\scanpage{4}\mspage{3}

Si $\deleted{\unsure{\delta|g}}\lambda$ est une section de $\mathrm{NS}_{\deleted{P}B/S}$ [donc définit une section, notée encore $\lambda$, de $\underline{\mathrm{NS}}_{P^\delta/S}$], alors $\underline{\mathrm{Pic}}^\lambda_{P^\delta/S}$ se déduit de $\underline{\mathrm{Pic}}^\lambda_{B/S}$ en tordant à l'aide de $P^\delta$. Or soit

\[
A=\underline{\mathrm{Pic}}^{\unsure{\circ|0}}_{B/S}\quad(=\underline{\mathrm{Alb}}^\circ_{X/S})
\]

le schéma abélien dual de $B$, \deleted{\unsure{$(=\underline{\mathrm{Alb}}^\circ($}} alors $\lambda$ définit un hom

\[
\widetilde\lambda:B\longrightarrow A
\]

et l'action de $B$ sur $\underline{\mathrm{Pic}}^\lambda_{B/S}$ se déduit de l'action évidente de $A=\underline{\mathrm{Pic}}^\circ_{B/S}$ sur $\underline{\mathrm{Pic}}^\lambda_{B/S}$ (qui est un $A$-torseur) via l'hom $\widetilde\lambda$. Donc on trouve

\[
\underline{\mathrm{Pic}}^\lambda_{P^\delta/S}\simeq\underline{\mathrm{Pic}}^\lambda_{B/S}
\overset{A}{\times}
\underbrace{\left(P^\delta\overset{B}{\times} A\right)}_{
\substack{A\text{-torseur déduit}\\\text{du }B\text{-torseur}\\ P^\delta\text{ grâce à l'hom.}\\ \widetilde\lambda:B\to A}}
\]

\deleted{\unsure{No|Alor|Alo}.} Reprenant les notations de c), nous allons appliquer ceci au cas où l'on a

\[
\lambda=\varphi(\delta),
\]

de sorte que, pour \deleted{\unsure{la|tt}} section $g$ de $X/S$, la restriction $L_g\mid P^\delta$ définit une section de $\underline{\mathrm{Pic}}^\lambda_{P^\delta/S}$, d'où une section
